\documentclass{article}
\usepackage{amsmath}
\begin{document}
\title{Submitting finished work can change which unfinished work remains feasible}
\maketitle

\paragraph{Source and candidate seam.}
In \texttt{usnistgov/ARIAC} tag 2024.3.0, the selected valid
\texttt{SubmitOrderServiceCallback} sets the named result, calls the
scorer and records submission membership. Separate selected native
branches of the world update test that same membership for
dependent-order announcement, robot malfunction and sensor blackout.
The order in the source update is to process announcements, process
challenge starts, process expiry, then publish health when due.
The selected replay preserves the relevant start order and simulates
neither expiry nor asynchronous delivery. The scorer and publication
bodies are stubs: measured outputs are calls and local state changes,
not scores or robot motion. Each history uses a valid preannounced
order and one unchanged trial.

\paragraph{Composition failure.}
Acknowledging the submission and choosing the time at which a
robot must tolerate a new fault are different operations, even in
the same simulator. Two serialized histories can have the very same
old health flags before the world update. One already has a
committed old order with a lost result. The other has not yet
committed it. The next update can create the new duty and revoke the
availability of the floor robot and logical camera in the first
history, without any changed pre-update health flag. Valid submission
after that update changes the challenge cut by a full modeled update
period. If retried after a lost response inside the same trial, the
original submit service rejects the repetition, without running the
scorer or announcing the descendant a second time. That rejection
does not transport the original result. This is a same-trial control,
not a result about server restart, message loss probabilities, or
unknown order membership.

\paragraph{Exact staged-cut check in the scoped family.}
Suppose source update ticks are consecutive integers.
A selected valid submission occurs just after tick \(c\).
A descendant must be announced by tick \(D\). The old floor duty
needs the floor's availability through its effect at \(f\); the
old camera duty likewise needs uninterrupted source availability
through tick \(s\). These are authored effect requirements, not
observed robot tasks. Let \(F\subseteq\{f,s\}\) contain the ends
of just those duties on resources disabled by the submission's
configured challenge. Other continuing work is independent in
this fixture. Disablements last beyond the test horizon.
The law for an admissible submit is
\[
  \max(\{0\}\cup F)\le c,\qquad c+1\le D.
\]
The first inequality protects the old effects against the first
possible loss tick \(c+1\); the second meets the descendant
deadline. Choosing the earliest such cut is exact. The replay
compares all eight opportunity choices, both challenge kinds
alone and together, and deliberate outside-envelope deadline
conditions. It executes native start flags and announcement for
each choice. A direct scheduler given identical ongoing duties,
authority, and transition rules agrees; it can encode these
as ordinary conditioned temporal effects. The gain over an eager
submission does not establish superiority to that baseline.

\paragraph{Where it stops.}
Unlike a fixed future opportunity, here the act can create the
opportunity, a new duty, and the future disturbance together.
Treating a pending submission as if the caller could still choose
its cut is not justified. A transport-level retry cannot remove
a committed simulator effect. For unconfirmed old submissions
a stronger gate would need authoritative commit identity or an
independent effect-side way to protect ongoing work. The fixture
does not show control over an already committed submission, an
adapter for real co-executing hardware, or a guarantee for a
second conditional fault whose onset is hidden. It is development
evidence for a specific prior-work collision. No independent case
has been chosen or used.

\paragraph{Relation to the nearest named prior architecture.}
The HeRo ICRA 2025 abstract and author poster describe state-machine
coordination with periodically synchronized state, monitoring
around execution, and planning-generated recovery in ARIAC.
The poster gives this workflow-level context, not the full solution
or the source code. This replay establishes a distinction any
comparison using the pinned ARIAC branches must represent. It does
not establish that HeRo omitted it, used this exact release,
would fail these tests, or is reconstructible from the poster.
See DOI \texttt{10.1109/ICRA55743.2025.11128674}.
\end{document}
